\documentclass[10pt,a4paper]{amsart}
\usepackage{amsmath,amssymb,amsthm,graphicx}
\usepackage{verbatim}
\newenvironment{reference}{\comment}{\endcomment}
%\newenvironment{reference}{}{}
\newenvironment{slogan}{\comment}{\endcomment}
\newenvironment{history}{\comment}{\endcomment}
% For commutative diagrams we use Xy-pic
\usepackage[all]{xy}

% We use 2cell for 2-commutative diagrams.
\xyoption{2cell}
% Theorem environments.
%
\theoremstyle{plain}
\newtheorem{theorem}[subsection]{Theorem}
\newtheorem{proposition}[subsection]{Proposition}
\newtheorem{lemma}[subsection]{Lemma}

\theoremstyle{definition}
\newtheorem{definition}[subsection]{Definition}
\newtheorem{example}[subsection]{Example}
\newtheorem{exercise}[subsection]{Exercise}
\newtheorem{situation}[subsection]{Situation}

\theoremstyle{remark}
\newtheorem{remark}[subsection]{Remark}
\newtheorem{remarks}[subsection]{Remarks}

\numberwithin{equation}{subsection}

% Macros
%
\def\lim{\mathop{\mathrm{lim}}\nolimits}
\def\colim{\mathop{\mathrm{colim}}\nolimits}
\def\Spec{\mathop{\mathrm{Spec}}}
\def\Hom{\mathop{\mathrm{Hom}}\nolimits}
\def\Ext{\mathop{\mathrm{Ext}}\nolimits}
\def\SheafHom{\mathop{\mathcal{H}\!\mathit{om}}\nolimits}
\def\SheafExt{\mathop{\mathcal{E}\!\mathit{xt}}\nolimits}
\def\Sch{\mathit{Sch}}
\def\Mor{\mathop{\mathrm{Mor}}\nolimits}
\def\Ob{\mathop{\mathrm{Ob}}\nolimits}
\def\Sh{\mathop{\mathit{Sh}}\nolimits}
\def\NL{\mathop{N\!L}\nolimits}
\def\CH{\mathop{\mathrm{CH}}\nolimits}
\def\proetale{{pro\text{-}\acute{e}tale}}
\def\etale{{\acute{e}tale}}
\def\QCoh{\mathit{QCoh}}
\def\Ker{\mathop{\mathrm{Ker}}}
\def\Im{\mathop{\mathrm{Im}}}
\def\Coker{\mathop{\mathrm{Coker}}}
\def\Coim{\mathop{\mathrm{Coim}}}

% Boxtimes
%
\DeclareMathSymbol{\boxtimes}{\mathbin}{AMSa}{"02}

%
% Macros for moduli stacks/spaces
%
\def\QCohstack{\mathcal{QC}\!\mathit{oh}}
\def\Cohstack{\mathcal{C}\!\mathit{oh}}
\def\Spacesstack{\mathcal{S}\!\mathit{paces}}
\def\Quotfunctor{\mathrm{Quot}}
\def\Hilbfunctor{\mathrm{Hilb}}
\def\Curvesstack{\mathcal{C}\!\mathit{urves}}
\def\Polarizedstack{\mathcal{P}\!\mathit{olarized}}
\def\Complexesstack{\mathcal{C}\!\mathit{omplexes}}
% \Pic is the operator that assigns to X its picard group, usage \Pic(X)
% \Picardstack_{X/B} denotes the Picard stack of X over B
% \Picardfunctor_{X/B} denotes the Picard functor of X over B
\def\Pic{\mathop{\mathrm{Pic}}\nolimits}
\def\Picardstack{\mathcal{P}\!\mathit{ic}}
\def\Picardfunctor{\mathrm{Pic}}
\def\Deformationcategory{\mathcal{D}\!\mathit{ef}}
\usepackage{cleveref}
\crefname{originallemma}{original lemma}{original lemmas}
\Crefname{originallemma}{Original lemma}{Original lemmas}
\makeatletter
\newlabel{coherent-lemma-affine-morphism-projection-ideal}{{01YP}{source}{}{sourcepointer-1}{}}
\newlabel{coherent-lemma-affine-morphism-projection-ideal@cref}{{[originallemma][0][]01YP}{[1][0][]source}{}{}{}}
\makeatother
\title{The ideal of the target acting on a pushforward along an affine morphism, with the notation spelled out}
\author{The Stacks Project Authors}
\date{stacks-project:spaces-cohomology.tex:local-snapshot-20261012; GNU FDL 1.2 or later}
\begin{document}
\maketitle
\noindent\textit{Editorial note.} All delivered mathematical text is the author's own, copied from the pinned Stacks Project chapter file spaces-cohomology.tex; no statement or proof step was rewritten, reordered or removed.\par\medskip
\noindent\textit{Editorial note.} Preamble provenance: the pinned chapter file opens with \textbackslash{}input\{preamble\} and carries no preamble text of its own, so the project's single Stacks-wide preamble staged at sources/stacks-spaces-cohomology/preamble.tex is excerpted verbatim as three spans, lines 8-12 (verbatim package and the project's own reference/slogan/history display environments), lines 14-18 (Xy-pic setup) and lines 154-217 (theorem declarations and mathematical macros). The wrapper class and the amsmath/amssymb/amsthm/graphicx base packages are supplied by the assembler. No macro was invented or redefined.\par\medskip
\noindent\textit{Editorial note.} Same-chapter pointers are located at their \textbackslash{}label lines inside the pinned chapter file; cross-chapter pointers are located in the project's own tag list staged in the snapshot (sources/stacks-chow/tags.txt), at the verbatim lines carrying their canonical tags.\par\medskip
\noindent\textit{Allowed standard prerequisites (editorial).}\begin{itemize}
\item The scheme case of the projection formula for an ideal along an affine morphism. Cohomology of Schemes, label coherent-lemma-affine-morphism-projection-ideal (invoked at proof line 3409). After the reduction to an affine base it is exactly the asserted equality \$\textbackslash{}mathcal\{I\}f\_*\textbackslash{}mathcal\{F\} = f\_*(f\textasciicircum{}\{-1\}\textbackslash{}mathcal\{I\}\textbackslash{}mathcal\{F\})\$.
\item Exactness of the inverse image functor, which makes the inverse image of an ideal a sheaf of ideals. Standard, used at proof lines 3396-3401. It shows that \$f\textasciicircum{}\{-1\}\textbackslash{}mathcal\{I\}\$ is a sheaf of ideals of \$f\textasciicircum{}\{-1\}\textbackslash{}mathcal\{O\}\_X\$, acting on \$\textbackslash{}mathcal\{F\}\$ through \$f\textasciicircum{}\textbackslash{}sharp\$, so that \$f\textasciicircum{}\{-1\}\textbackslash{}mathcal\{I\}\textbackslash{}mathcal\{F\}\$ makes sense.
\end{itemize}\medskip
\ifdefined\hypertarget\hypertarget{sourcepointer-1}{}\fi
\noindent\textit{Source reference.} Original lemma 01YP: sources/stacks-chow/tags.txt, source lines 2564--2564. Navigation pointer for \textbackslash{}ref\{coherent-lemma-affine-morphism-projection-ideal\} used inside the counted span (line 3413); the argument that consumes it is retained in the span.\par\medskip
\begin{lemma}
\label{lemma-affine-morphism-projection-ideal}
Let $S$ be a scheme. Let $f : Y \to X$ be a morphism of algebraic spaces
over $S$. Let $\mathcal{F}$ be a quasi-coherent sheaf on $Y$.
Let $\mathcal{I}$ be a quasi-coherent sheaf of ideals on $X$.
If $f$ is affine then
$\mathcal{I}f_*\mathcal{F} = f_*(f^{-1}\mathcal{I}\mathcal{F})$
(with notation as explained in the proof).
\end{lemma}

\begin{proof}
The notation means the following. Since $f^{-1}$ is an exact functor
we see that $f^{-1}\mathcal{I}$ is a sheaf
of ideals of $f^{-1}\mathcal{O}_X$. Via the map
$f^\sharp : f^{-1}\mathcal{O}_X \to \mathcal{O}_Y$
on $Y_\etale$ this acts on $\mathcal{F}$.
Then $f^{-1}\mathcal{I}\mathcal{F}$ is the subsheaf
generated by sums of local sections of the form $as$ where $a$
is a local section of $f^{-1}\mathcal{I}$ and $s$ is a local section
of $\mathcal{F}$. It is a quasi-coherent $\mathcal{O}_Y$-submodule
of $\mathcal{F}$ because it is also the image of a natural map
$f^*\mathcal{I} \otimes_{\mathcal{O}_Y} \mathcal{F} \to \mathcal{F}$.

\medskip\noindent
Having said this the proof is straightforward. Namely, the question is
\'etale local on $X$ and hence we may assume $X$ is an affine scheme.
In this case the result is a consequence of the corresponding result
for schemes, see Cohomology of Schemes, Lemma
\ref{coherent-lemma-affine-morphism-projection-ideal}.
\end{proof}

\end{document}
